% !TeX root = ../../main.tex
\section{算符函数及其相关概念}

\begin{solution}[算符函数与幂级数]
\problem{我们知道，幂级数
\[
f(x)=\sum_{n=0}^{\infty}x^n,
\]
在 $|x|<1$ 时等价于函数 $f(x)=(1-x)^{-1}$。现考虑通过过渡到本征基，检验一个厄米算符 $\Omega$ 的 $q$ 数幂级数
\[
f(\Omega)=\sum_{n=0}^{\infty}\Omega^n,
\]
在何种条件下可以被视为 $(1-\Omega)^{-1}$.}

\answer
首先回顾标量情形：当 $|x|<1$ 时几何级数收敛并满足
\[
\sum_{n=0}^{\infty}x^n=(1-x)^{-1}.
\]

对于算符情形，设 $\Omega$ 为厄米算符，根据谱定理可在某个正交归一基下对角化：
\[
\Omega=\mathrm{diag}(\omega_1,\omega_2,\cdots,\omega_N),
\]
其中 $\omega_i\in\mathbb{R}$ 为本征值。

在该本征基中，幂级数逐项作用为
\[
\Omega^n=\mathrm{diag}(\omega_1^n,\omega_2^n,\cdots,\omega_N^n),
\]
因此
\[
f(\Omega)=\sum_{n=0}^{\infty}\Omega^n
=\mathrm{diag}\left(\sum_{n=0}^{\infty}\omega_1^n,\ \sum_{n=0}^{\infty}\omega_2^n,\ \cdots\right).
\]

只要每个本征值满足几何级数收敛条件 $|\omega_i|<1$，则逐个分量求和得到
\[
\sum_{n=0}^{\infty}\omega_i^n=\frac{1}{1-\omega_i}.
\]
于是
\[
f(\Omega)=\mathrm{diag}\left(\frac{1}{1-\omega_1},\frac{1}{1-\omega_2},\cdots\right).
\]

另一方面，
\[
(1-\Omega)^{-1}
=\mathrm{diag}\left(\frac{1}{1-\omega_1},\frac{1}{1-\omega_2},\cdots\right).
\]

因此在谱半径满足
\[
\max_i |\omega_i|<1
\]
的条件下，幂级数收敛且成立算符恒等式
\[
\sum_{n=0}^{\infty}\Omega^n=(1-\Omega)^{-1}.
\]

结论：当且仅当 $\Omega$ 的全部本征值满足 $|\omega_i|<1$（即谱半径小于1）时，该算符幂级数收敛，并可严格视为 $(1-\Omega)^{-1}$。
\end{solution}


\begin{solution}[算符指数的幺正性]
\problem{如果 $H$ 是一个厄米算符，证明 $U=e^{iH}$ 是幺正的。（请注意与 $c$ 数的类比：如果 $\theta$ 是实数，则 $u=e^{i\theta}$ 是一个模为 1 的数。）}

\answer
目标是证明 $U=e^{iH}$ 满足幺正性条件 $U^\dagger U=I$。

首先利用厄米算符的基本性质：$H^\dagger=H$。考虑算符指数的厄米共轭。由幂级数定义
\[
e^{iH}=\sum_{n=0}^\infty \frac{(iH)^n}{n!},
\]
对其取厄米共轭并利用 $(AB)^\dagger=B^\dagger A^\dagger$ 以及 $H^\dagger=H$ 得
\[
U^\dagger=(e^{iH})^\dagger=\sum_{n=0}^\infty \frac{(-iH)^n}{n!}=e^{-iH}.
\]

因此
\[
U^\dagger U = e^{-iH}e^{iH}.
\]

由于 $H$ 与自身显然对易，指数算符满足合并规律
\[
e^{-iH}e^{iH}=e^{(-iH+iH)}=e^0=I.
\]

因此得到
\[
U^\dagger U=I,
\]
同理可得 $UU^\dagger=I$，故 $U$ 为幺正算符。

另一种等价观点是利用谱分解：由于 $H$ 为厄米算符，可写为
\[
H=\sum_n \omega_n |n\rangle\langle n|,
\]
则
\[
U=e^{iH}=\sum_n e^{i\omega_n}|n\rangle\langle n|.
\]
其厄米共轭为
\[
U^\dagger=\sum_n e^{-i\omega_n}|n\rangle\langle n|,
\]
逐本征子相乘即可得到单位算符，从而再次证明幺正性成立。

结论：当 $H$ 为厄米算符时，$U=e^{iH}$ 必为幺正算符。
\end{solution}


\begin{solution}[行列式与算符函数]
\problem{对于上述情况，证明 $\det U = e^{i\,\mathrm{Tr}\,H}$.}

\answer
设 $H$ 为厄米算符，根据谱定理可在某一正交归一基下对角化为
\[
H=\mathrm{diag}(\omega_1,\omega_2,\cdots,\omega_n),\quad \omega_i\in\mathbb{R}.
\]
由算符指数的定义，$U=e^{iH}$ 在同一基下亦为对角矩阵，其对角元逐一为
\[
U=\mathrm{diag}\left(e^{i\omega_1},e^{i\omega_2},\cdots,e^{i\omega_n}\right).
\]

因此行列式可直接取对角元乘积：
\[
\det U=\prod_{k=1}^n e^{i\omega_k}.
\]
利用指数性质得到
\[
\det U=\exp\left(i\sum_{k=1}^n \omega_k\right).
\]

另一方面，由于迹在对角基下等于本征值之和，
\[
\mathrm{Tr}\,H=\sum_{k=1}^n \omega_k,
\]
因此可得
\[
\det U=\exp\left(i\,\mathrm{Tr}\,H\right).
\]

若用基底无关表述，上述结论可理解为：行列式与迹均在相似变换下不变，因此在任意表示中成立。
\end{solution}